\section{Reproductibilité et reprise du code}\label{sec:repro-ref}

\subsection{Environnement}
\begin{itemize}
\item \textbf{Langages et outils.} Python, géré par \texttt{uv} (\texttt{pyproject.toml}, \texttt{uv.lock}) ;
  Node.js pour la recette de la carte et la vidéo ; \texttt{tectonic} pour les articles.
\item \textbf{Dépôt.} \texttt{progdinh/hackathon\_acces\_soins}, dossier \texttt{v2/}. Le code utile de la v1 y
  est absorbé (\texttt{v2/sante/}) ; ses registres sont archivés en lecture seule dans
  \texttt{v2/docs/lab/archive-v1/}.
\item \textbf{Données.} Les 34 fichiers traités sont commités. Chacun est listé dans
  \texttt{data/processed/MANIFEST.md} et \texttt{v2/data/processed/MANIFEST.md}, avec sa taille, son empreinte
  SHA-256, son script, sa commande et sa source. Un clone neuf suffit pour lancer la carte. \texttt{jumeau.check} demande en plus un
  fichier brut non commité, ouvert : le jeu DREES 491 (à télécharger ;
  \texttt{data/raw/drees/491\_formation\_series\_longues\_1966\_2025.xlsx}, lu par C16). Jusqu'au 15 novembre 2026,
  C5, C9 et la variante filière lisaient aussi un fichier d'inscrits fourni à l'équipe sous accord de
  confidentialité ; ils lisent désormais SISE (\texttt{v2/data/build\_sise.py}, Licence Ouverte), et
  \texttt{v2/data/manifest.py --check} échoue si un code exécuté par l'outil lit un fichier des organisateurs
  (liste \texttt{INTERDITS}, dossiers \texttt{EXECUTE}). Seul \texttt{v2/methodes/c5.py}, script exploratoire de
  l'axe~2 appelé par aucun autre code, lit encore ce fichier. Pour le reste, les données brutes (environ 3,6\,Go\mesure{} au total) ne servent qu'à reconstruire les
  tables.
\end{itemize}

\subsection{Commandes}
Toutes les commandes se lancent depuis la racine du dépôt.
\begin{small}
\begin{verbatim}
uv sync
uv run python v2/data/manifest.py --check          # manifestes à jour, aucun fichier
                                                   # des organisateurs dans la construction
PYTHONPATH=v2 uv run python -m jumeau.check        # 18 critères ; exige le jeu ouvert DREES 491
uv run python v2/prototype-carte/server.py         # carte, http://localhost:8767

# axe 2 : banc des méthodes et contrôles du 15 novembre
PYTHONPATH=v2 uv run python -m methodes.banc
PYTHONPATH=v2 uv run --with scipy python -W ignore -m methodes.banc_cohortes
PYTHONPATH=v2 uv run --with scipy python -W ignore -m methodes.budget_erreur
PYTHONPATH=v2 uv run --with scipy python -W ignore -m methodes.c1_intervalles tirages --level dep
PYTHONPATH=v2 uv run --with scikit-learn==1.9.1 python -m methodes.c1_gbm

# axe 3 : validation
PYTHONPATH=v2 uv run --with scipy python -m validation.did
PYTHONPATH=v2 uv run --with scipy python -m validation.dd76

# compilateur de langage naturel, jeu à l'aveugle n° 6
uv run python v2/langage/evaluer_hors_echantillon.py cas_test_hors_echantillon_6.yaml

# reconstruction d'une table depuis les données brutes (exemples)
PYTHONPATH=v2 uv run python -m sante.load
uv run python v2/data/build_secteur.py
uv run python v2/data/manifest.py                  # réécrit les manifestes
\end{verbatim}
\end{small}
\texttt{jumeau.check} écrit \texttt{v2/artifacts/a0\_jumeau/results.json}, qui n'est pas versionné et se
régénère à chaque run. Les sorties des scripts de l'axe 2 sont écrites dans \texttt{v2/methodes/out/}, celles de l'axe 3 dans
\texttt{v2/validation/out/}. La carte et l'assistant ont leurs propres commandes, données dans
\texttt{README.md} et \texttt{chatbot/README.md}. L'assistant demande une clé \texttt{ANTHROPIC\_API\_KEY} dans un
fichier \texttt{.env} ignoré par git ; sans clé, il bascule en mode déterministe.

\subsection{Où se trouvent les choses}
\begin{center}\small
\begin{tabular}{@{}ll@{}}
\toprule
Contenu & Chemin \\
\midrule
Jumeau (état, modules, scénarios, critères) & \texttt{v2/jumeau/} \\
Estimateurs, bancs, contrôles & \texttt{v2/methodes/} \\
Protocoles pré-enregistrés, dépôts, évaluations & \texttt{v2/validation/} \\
Construction des tables & \texttt{v2/data/build\_*.py}, \texttt{v2/sante/} \\
Compilateur de langage naturel & \texttt{v2/langage/} \\
Serveur et API & \texttt{v2/prototype/server.py} \\
Carte & \texttt{v2/prototype-carte/} \\
Assistant & \texttt{chatbot/}, \texttt{v2/agent/} \\
Articles par axe et ce rapport & \texttt{v2/papers/} \\
Registre des décisions & \texttt{v2/docs/lab/decisions.md} \\
Questions ouvertes & \texttt{v2/docs/lab/open.md} \\
Rapports, dérivations, fiches de paramètres & \texttt{v2/docs/lab/\{reports,derivations,fiches\}/} \\
\bottomrule
\end{tabular}
\end{center}

\subsection{Règles de travail à conserver}
Ces règles ont structuré le projet ; elles sont reprises de \texttt{CLAUDE.md} et \texttt{v2/ROADMAP.md}.
\begin{itemize}
\item Tout estimateur passe par l'interface commune, et toute affirmation de précision est comparée à une baseline
  naïve sur le même banc.
\item Le test qui peut échouer est écrit avant le calcul. Le dépôt de ce test devra être horodaté hors de git
  (section~\ref{sec:validation-ref}).
\item Une décision se note au registre, datée. Une décision remplacée reste visible, marquée comme remplacée.
\item Aucune donnée brute ni aucun fichier sous accord de confidentialité n'entre dans le dépôt.
\end{itemize}
